\section{Most recent impactful developments}
This section prioritizes recent work with direct implications for our reward-tube project.
The ordering reflects relevance to our research question rather than a bibliometric
ranking; recent preprints are distinguished from conference publications in their entries.

\subsection{Research question and shared notation}
\label{sec:review}
\textbf{Question.} Can a calibrated, structured uncertainty set around a process
reward model (PRM) reduce reward hacking under inference-time search or
reinforcement learning, while retaining useful discrimination between reasoning paths?
We compare with existing pessimistic methods at matched compute, annotation budget,
aggregation, and policy movement.

This review was expanded on October 7, 2026. Each entry gives the mathematical
object being modeled, important conditions or empirical limitations, and a concrete
connection to our project. Equations use consistent notation rather than always
reproducing the paper's symbols. Source section/equation pointers refer to the
specified version where given; numbering may change in later revisions. A
\emph{schematic} objective describes a method's structure rather than an exact
implementation. Statements introduced as \emph{our interpretation}, \emph{our
diagnostic}, or \emph{our extension} are research synthesis, not the authors' results.
We inspected relevant method and theory sections; this is not a complete proof audit
or a reproduction of the empirical findings.

\subsubsection{Shared notation}
Let $x$ be a problem, $a_t$ a reasoning step, $h_t=(x,a_{1:t})$ its prefix,
and $\tau=(a_1,\ldots,a_T)$ a trajectory. Write $\hat r(h_t)$ for a learned
process score and $G(\hat{\bm r}(\tau))$ for its trajectory aggregation.
Intrinsic step validity, validity of the maintained reasoning state, and
\[
 p_{\pi_c}(h_t)=\Prob_{\pi_c}\{\text{correct final answer}\mid h_t\}
\]
are distinct targets; $\pi_c$ is a specified continuation policy. For additive
rewards, $J_r(\pi)$ denotes expected return. When using normalized discounted
occupancies, $d_\pi$ is a probability distribution and $J_r(\pi)=\E_{d_\pi}r$.
Set $\Delta J_r=J_r(\pi)-J_r(\pi_0)$ for a reference policy $\pi_0$.
Two recurring losses are
\[
 \ell_{\rm BCE}(y,p)=-y\log p-(1-y)\log(1-p),\qquad
 \rho_\beta(u)=u\bigl(\beta-\ind\{u<0\}\bigr).
\]
The second is the pinball loss for a $\beta$-quantile. A prediction interval
covers a future observed label; a function confidence set aims to contain one
unknown reward function. Neither term should substitute for the other.

\subsection{Uncertainty modeling and pessimistic reward optimization}

\paper{Ma et al., Distributional PRMs via Conditional Optimal Transport (2026; arXiv v1)}{
\textbf{Model.} PRM features condition a one-dimensional OT map. Convex
potentials parameterized by partially input-convex networks give
\[
 T_\theta(u,h)=\partial_u g_\theta(u,h),\qquad
 T_\theta(\cdot,h)_\#\mu_h\approx\nu_h.
\]
Here $\nu_h$ models rollout-success targets. The conditional dual optimization
trains the potentials; a resulting quantile function $Q_\theta(\beta\mid h)$
can be queried at arbitrary levels \cite{otprm}.

\textbf{Mathematical detail.} Convexity makes the transport map nondecreasing,
avoiding quantile crossing. For a general continuous source CDF $F_{\mu_h}$,
standard OT quantile extraction is
$Q_{\nu_h}(\beta)=T(F_{\mu_h}^{-1}(\beta),h)$; $T(\beta,h)$ itself is a
quantile only with the corresponding uniform-source parameterization.
Monotonicity is a structural property, not a finite-sample coverage theorem.

\textbf{Evidence and connection.} The paper reports calibration and adaptive
budget improvements using empirical completion-success labels. Keep its
calibration/evaluation split explicit. Our extension needs joint process
uncertainty, calibrated model discrepancy, or post-selection assurance; learning
a flexible distribution or monotone quantiles is already this paper's contribution.
Changing the completer changes its target.}

\paper{Yu et al., From Curiosity to Caution (ICLR 2026; arXiv v1)}{
\textbf{Score and learning.} Random-network distillation fits predictor
$P_\theta$ to fixed target features $T$ on reference responses. At inference,
\[
 \alpha(x,y)=\norm{P_\theta(x,y)-T(x,y)}^2,
 \quad r_{\rm caution}=\hat r-\lambda\alpha,
 \quad \hat i=\arg\max_{i\leq N}r_{\rm caution}(x,y_i).
\]
Section 2.3 gives a stylized linear-reward/RND analysis whose informal theorem
states a reward gain of order $\sqrt{\log N}$ over ordinary BoN under the stated
sampling and network conditions \cite{caution}.

\textbf{Conditions and evidence.} The theory is a proof of concept under
strong distribution/feature assumptions, not a theorem for arbitrary LLMs.
The penalty measures reference novelty: unusual valid reasoning may be penalized,
and familiar invalid reasoning can evade it. Predictor training needs unlabelled
reference responses; scoring requires extra network passes.

\textbf{Connection.} An important low-cost uncertainty comparator. Check whether
calibrated residual uncertainty improves on novelty alone, at matched generation
and scoring cost. The name ``LCB'' does not establish calibrated reward-error
coverage in the actual PRM setting.}

\paper{Hahami et al., A Unifying Lens on Reward Uncertainty (2026; arXiv v1)}{
\textbf{Variational formula.} For believed reward distribution $P$, a
KL-penalized reward adversary gives
\[
 r_{\rm eff}=\inf_Q\{\E_Q r+\eta\KL(Q\|P)\}
 =-\eta\log\E_P e^{-r/\eta}.
\]
The optimizer exponentially tilts $P$ toward lower rewards. For a uniform
ensemble, $r_{\rm eff}=-\eta\log[K^{-1}\sum_k e^{-r_k/\eta}]$; see Section 3
\cite{unifying}.

\textbf{Limits and assumptions.} Under suitable moments,
$r_{\rm eff}=\E r-\Var(r)/(2\eta)+O(\eta^{-2})$ as $\eta\to\infty$;
$\eta\downarrow0$ gives the essential minimum. The Gaussian expression is
exact when $r$ really is Gaussian. Exponential moments must exist. Reward
adversary temperature $\eta$ and policy KL coefficient have different roles,
even if a particular formulation sets them equal. A hard-radius KL ball instead
requires optimizing a dual multiplier.

\textbf{Connection.} Mean, minimum, and variance penalties already share an
entropic formulation. Our extension should model a \emph{joint} step-reward
law and its calibration, rather than rediscover scalar log-sum-exp pessimism.
An incorrect believed distribution remains incorrect after robust tilting.}

\paper{Zhou et al., PRISM (ACL 2026)}{
\textbf{Model and likelihood.} Frozen features feed Gaussian expert heads and
a prompt-dependent router:
\[
 p(r\mid x,y)=\sum_{k=1}^{K}\omega_k(x)
 \mathcal N(r;\mu_k(x,y),\sigma_k^2(x,y)).
\]
Stage 1 forms experts; Stage 2 freezes them and learns router weights. Its
cross-expert preference likelihood averages
$\Phi((\mu_k(y_+)-\mu_l(y_-))/\sqrt{\sigma_k^2(y_+)+\sigma_l^2(y_-)+\epsilon})$
with weights $\omega_k\omega_l$, where $\Phi$ is the normal CDF; Section 3,
Eqs. (8)--(10) \cite{prism}.

\textbf{Conditions and interpretation.} Mixture structure represents preference
heterogeneity and within-component dispersion. These components are not
identified epistemic versus aleatoric uncertainty solely by their names; latent
mixtures and approximate likelihoods need empirical validation. Gaussian
reward differences use the model's independence convention.

\textbf{Connection.} A rich-distribution comparator for preference rewards,
not a certified process-verification method. Our extension could share latent
expert weights across steps, creating coherent trajectory-level scenarios.
Decide whether uncertainty concerns conflicting legitimate preferences or
incorrect mathematical judgments; averaging them can obscure the target of
the tube.}

\paper{Duan et al., Bayesian Non-negative Reward Modeling / BNRM (2026; arXiv v1)}{
\textbf{Generative model.} Local activations $z\geq0$ and a shared dictionary
$\phi\geq0$ have Gamma priors, generate reward $r=z^\top\phi$, and define
preference probability $\sigmoid(r_+-r_-)$. Weibull variational posteriors are
fit using the regularized evidence objective
\[
 \E_q\log p(\mathcal D\mid z,\phi)
 -\eta\KL(q(z)\|p(z))-\eta\KL(q(\phi)\|p(\phi)).
\]
See Section 4, Eqs. (5)--(9) \cite{bnrm}.

\textbf{Conditions and evidence.} Nonnegative sparse factors encourage
parts-based representations; they do not prove causal disentanglement or
elimination of style bias. A variational posterior can underestimate uncertainty,
and misspecification can persist across all factors. Preference rewards are
identified through differences, so positivity alone is not a correctness
certificate. Reported robustness gains concern the evaluated RLHF settings.

\textbf{Connection.} A possible alternative to Gaussian heads or independent
ensembles. A shared dictionary creates a natural source of correlated error
across prefixes; local latent randomness and global uncertainty should be
separated when constructing trajectory scenarios. Any posterior credible tube
still needs independent coverage and optimization stress tests.}

\subsection{Reward hacking diagnostics and adversarial defenses}

\paper{Tiwari et al., Reward Under Attack (2026; arXiv v1)}{
\textbf{Setup and attack objective.} The paper uses model-specific aggregation:
Skywork's terminal score and Qwen's minimum step score. Adversarial probing seeks
high-scoring token modifications; in our notation its attack target is
\[
 \max_{u\in\mathcal A}\ G(\hat{\bm r}(\tau\mathbin{\oplus}u)),
\]
subject to the attack's permitted placements and token budget. Static
perturbations, gradient-based token optimization, and GRPO expose different
failure mechanisms \cite{attack}.

\textbf{Evidence and limits.} In the studied Skywork/AIME setting, reward
approaches one while answer accuracy remains below 4\%. Qwen resists the
particular gradient attack but collapses during RL to a vacuous introductory
sentence. Therefore success against one attack is not general robustness.
PRM-BiasBench tests controlled content and presentation changes.

\textbf{Connection.} Optimize the \emph{deployed robust score}, not only the
nominal PRM. Our diagnostic should report reward, independent correctness,
claim count, length, and vacuity across pressure levels. Parser exploits and
reward inflation require separate tests; a score tube that covers only
mathematical steps may miss an output that avoids making any claim.}

\paper{Juneja et al., Adversarial Training for PRMs / APRM (ICML 2026)}{
\textbf{Game and objective.} A generator corrupts a reference step; a verifier
predicts correctness $\hat y$, checked by an algorithmic oracle $y$.
Verifier payoff is $r_R=2\ind\{y=\hat y\}-1$. Generator payoff is $+1$ for
an incorrect accepted step, $0$ for an incorrect rejected step, and $-1$ for
an uncorrupted correct step. Each player optimizes
\[
 U_i'=\E r_i-\tau\KL(\pi_i\|\pi_{i,0})+c_H H(\pi_i).
\]
The inspected preprint v1, Section 2, formulates a general-sum game
\cite{aprm}.

\textbf{Theoretical conditions.} Finite action spaces, continuity, and
concavity conditions support equilibrium arguments in mixed-policy space.
Own-player strong concavity is not by itself enough to make a general-sum
pseudo-gradient strongly monotone; cross-player derivatives matter. Neural PPO
optimization does not automatically inherit those policy-space properties.
The oracle uses reference-based semantic/structural checks and has its own error modes.

\textbf{Connection.} A strong hard-negative baseline and adversarial data source.
Compare the same annotation budget spent on retraining versus tube construction;
audit oracle mistakes before treating generated corruption labels as truth.}

\paper{Pathmanathan and Huang, REFORM (ACL 2026; inspected arXiv v1)}{
\textbf{Method and objective.} REFORM discovers preference-order reversals while
attempting to preserve response quality. An aligned policy proposes top-$k$
tokens and reward-guided decoding searches for low-scoring preferred responses;
a misaligned policy searches for high-scoring nonpreferred responses.
Failure-aware retraining minimizes Bradley--Terry loss,
\[
 \mathcal L_{\rm BT}=-\E\log\sigmoid\bigl(r_\phi(x,y_+)-r_\phi(x,y_-)\bigr),
\]
on original and discovered preference pairs \cite{reform}.

\textbf{Conditions and evidence.} Section 3 substitutes a prefix reward for
an intractable expected completion reward during controlled decoding. Policy
likelihood and known preference classes constrain the search, but high likelihood
does not certify class preservation. New examples still need reliable preference
labels. The v1 title differs from the final bibliography title.

\textbf{Connection.} Optimizer-discovered failures are useful calibration or
misspecification examples. Our extension should compare against retraining on
those same examples. Maintain an independent audit split: fitting both the
reward and its tube to discovered attacks can overfit the tested attack family.}

\subsection{New supervision targets and reasoning-state models}

\paper{Xie et al., Noise-aware Learning (2026; arXiv v1)}{
\textbf{Method and formula.} Reflection-aware correction rejects a successful
continuation as evidence for step $t$ if a later step explicitly repairs $t$.
Noise-aware iterative training then updates a filtered label via
\[
 \tilde y_t^{(k+1)}=
 \begin{cases}
 r_t^{(k)},&|r_t^{(k)}-\tilde y_t^{(k)}|>\delta,\\
 \tilde y_t^{(k)},&\text{otherwise},
 \end{cases}
\]
where $r_t^{(k)}$ is the current PRM prediction. This is the refinement rule in
Eq. (7), followed by retraining \cite{noise}.

\textbf{Conditions and evidence.} The correction targets MC false positives
caused by repair; limited completer capability can still produce false negatives.
Selective self-relabeling assumes that, after filtering, large disagreement is
more often label noise than model error. If that assumption fails, the procedure
can reinforce confident mistakes. Reported step-discrimination gains are
empirical, not finite-sample guarantees.

\textbf{Connection.} Treat reflection indicators and label/model disagreement
as candidate uncertainty features. Evaluate against independent human or formal
labels; the model's own refined labels cannot also serve as the truth used to
validate its tube.}

\paper{Wang et al., Hierarchical Multi-Step Reward Models (arXiv v5, 2026)}{
\textbf{Model.} HRM augments single-step supervision with consecutive merged
steps. Its basic dataset construction is
\[
 \mathcal D_{\rm HRM}=\mathcal D_{\rm PRM}\cup
 \{(a_i\mathbin{\oplus}a_{i+1},R(a_i\mathbin{\oplus}a_{i+1}))\}_{i<T},
\]
where $\oplus$ concatenates steps. Hierarchical node compression augments
MCTS-derived supervision. A two-class softmax produces process scores
\cite{hrm}.

\textbf{Conditions and evidence.} Section 3.2's compression argument preserves
a sample-mean expectation under its i.i.d.-child assumption, while changing
variance from $\sigma^2/N$ to $\sigma^2/(N-1)$ when one child is removed.
Correlated search children need separate treatment. The policy-improvement
experiments use filtered-data supervised fine-tuning rather than PPO/GRPO;
stronger search performance is not an RL robustness theorem.

\textbf{Connection.} Merged-step supervision is a representation baseline for
inter-step dependence. Our structured tube should improve on a well-trained
multi-step center, not merely on an impoverished independent-step model.
Ordinary prefix-conditioned PRMs already see history; explicit multi-step
supervision changes the training signal, not the existence of prefix context.}

\paper{Han et al., Cliff (September 2026; arXiv v1)}{
\textbf{Method and formula.} A teacher solves and verifies a problem, then
identifies the first pitfall in incorrect student rollouts. With binary group
success mean $\mu$ and $\sigma=\sqrt{\mu(1-\mu)}$, define
$A_{\rm cor}=(1-\mu)/\sigma$ and $A_{\rm inc}=-\mu/\sigma$. Cliff assigns
\[
 A_{ij}=
 \begin{cases}
 A_{\rm cor}-b,&\text{correct rollout},\\
 \lambda A_{\rm cor}-b,&\text{incorrect, before pitfall},\\
 A_{\rm inc}-b,&\text{incorrect, at/after pitfall}.
 \end{cases}
\]
Here $b$ is its centering baseline and $\lambda$ adjusts prefix credit; see
Section 3.2, Eqs. (3)--(4) \cite{cliff}.

\textbf{Conditions and evidence.} Incorrect teacher solutions trigger a
vanilla-GRPO fallback. Homogeneous outcome groups need the implementation's
zero-variance convention. Teacher validity and error localization are substantive
assumptions; a correct teacher answer does not certify every judgment.

\textbf{Connection.} Put uncertainty on the pitfall index rather than only
scalar step scores. Compare prefix/suffix credit with robust scoring. The first-error
boundary assumes a useful decomposition; repairing an earlier mistake can make
that decomposition too coarse.}

\paper{Gan et al., Reasoning State Propagation (September 2026; arXiv v1)}{
\textbf{Model and formula.} RSP predicts break probability $b_t$ and repair
probability $u_t$ for valid/invalid maintained reasoning states. With row-vector
state probabilities,
\[
 M_t=\begin{pmatrix}1-b_t&b_t\\u_t&1-u_t\end{pmatrix},\qquad
 \bm p_t=\bm p_{t-1}M_t,\qquad \bm p_0=(1,0).
\]
The process score is $p_t^{\rm valid}$. BCE on annotated intermediate states and
on the terminal outcome trains the propagated probabilities; Section 3,
Eqs. (2)--(9) \cite{rsp}.

\textbf{Conditions and evidence.} Outcome gradients reach earlier transitions
through the matrix product. Process labels ground intermediate states; outcome-only
data alone generally cannot identify every break/repair probability. Equating
a terminal valid state with answer correctness is a modeling decision that can
fail for lucky answers or incomplete reasoning.

\textbf{Connection.} A promising structured tube is a set of jointly admissible
transition matrices, not independent score intervals. Our extension could study
how transition uncertainty propagates and how ignoring repair causes excessive
pessimism. This requires a separate comparison with RSP and labels that distinguish
local errors from unresolved errors.}

\paper{Fan et al., TIPS (September 2026; arXiv v1)}{
\textbf{Method and theoretical detail.} A generative verifier emits thought $Z$,
step labels $\hat{\bm Y}$, and outcome judgment $\hat D$. GRPO uses only
$r=\ind\{\hat D=D\}$ to update the whole output. With pre-thought
representation $H$, outcome error $P_e$, and binary entropy $H_b$, the paper's
combined information bound is
\[
 I(Z;\bm Y\mid H)\geq
 H(D\mid H)-H_b(P_e)-H(D\mid\bm Y,H).
\]
This follows from its Theorems 1--2, Section 3.3 \cite{tips}.

\textbf{Conditions and evidence.} The bound is informative when verification
is nontrivial, outcome error is small, and process labels largely determine
the outcome. It concerns information in the thought, not correctness of each
emitted label, calibrated probabilities, or resistance to reward optimization.
Thought may encode the relevant information without the parser exposing it reliably.

\textbf{Connection.} A potentially inexpensive center or independent critic.
Audit generated labels separately and measure stochastic-judgment variability.
Our uncertainty set cannot inherit process validity from outcome supervision
without these additional checks.}

\subsection{Robust objectives and calibration theory}

\paper{Liu, Sun, and Zheng, Robust Optimization with Correlated Proxies (ICLR 2026)}{
\textbf{Uncertainty set and exact objective.} Normalize proxy $z$ to mean zero
and variance one under $d_{\pi_0}$. The paper constrains plausible rewards to
reference mean $M$, standard deviation $V>0$, and correlation $\rho$ with $z$.
Let $a=\E_{d_\pi}z$ and $D=\chi^2(d_\pi\|d_{\pi_0})$. Its inner solution gives
\[
 \inf_{r\in\U_{\rm corr}}J_r(\pi)
 =M+V\left[\rho a-\sqrt{1-\rho^2}\sqrt{D-a^2}\right].
\]
See the proceedings paper, Section 3.1, Eqs. (4)--(9) \cite{robustcorr}.

\textbf{Conditions and theory.} Support overlap and reference moments are
essential; outside reference support, the unconstrained adversary can make
rewards arbitrarily poor. The subtraction $D-a^2$ removes movement along the
normalized proxy direction. A known-feature linear variant restricts reward
structure and improves interpretability. The paper also analyzes occupancy
estimation and optimization.

\textbf{Connection.} A directly relevant example of \emph{coupling}, not an
independent interval construction. Our extension would estimate a process-informed
set from supervision and compare its geometry with this correlation set.
Generic max--min reward optimization and safe-reference improvement are
already prior work.}

\paper{Liu et al., Robust General Utility for RL (August 2026; arXiv v1)}{
\textbf{Objective.} The paper uses costs $f_\xi(\lambda_\theta)$ of policy
occupancy $\lambda_\theta$ and solves
\[
 \min_{\theta\in\Theta}\max_{\xi\in\Xi}f_\xi(\lambda_\theta),\qquad
 \Xi=\{\xi:d(\xi,\hat\xi)\leq d_0\}.
\]
Reward maximization reverses these signs. This includes linear rewards but also
nonlinear occupancy utilities; see Section 3, Eq. (2) \cite{generalutility}.

\textbf{Theoretical conditions.} The analysis separates utility regimes and
uses convex compact parameter sets, smoothness conditions, and appropriate
stationarity measures; proximal methods handle more general minimax cases.
A stationarity guarantee is not a global-optimality or true-reward guarantee
for arbitrary neural policies. The actual evaluation utility must belong to
the modeled ambiguity set for the robust interpretation to hold.

\textbf{Connection.} Useful when uncertainty concerns a utility beyond additive
step reward. However, trajectory minimum/product aggregation is not automatically
a function of ordinary state-action occupancy: history augmentation or a richer
trajectory distribution may be required. Our formulation should establish that
representation before importing general-utility algorithms or convergence claims.}

\paper{Conformal prediction through hypothesis testing \cite{conformaltest}}{
\textbf{Theoretical viewpoint.} Invert a family of level-$\alpha$ tests of a
candidate augmented observation:
\[
 C_n(x)=\{y:\varphi(Z_{1:n},(x,y))=0\}.
\]
If the test rejects the true augmented observation with probability at most
$\alpha$, then $\Prob\{Y_{n+1}\in C_n(X_{n+1})\}\ge1-\alpha$.
The paper studies universality of permutation-based constructions under
exchangeability, impossibility of unrestricted conditional coverage, and oracle
optimality. Its oracle construction uses inverse conditional-density scores.

\textbf{Conditions and limits.} Optimality is relative to a specified distribution
and efficiency criterion; the conditional density is ordinarily unknown. Marginal
exchangeability-based validity does not imply useful coverage at every rare input.
Continuous interval-size objectives also need adaptation for binary labels and
finite prediction sets.

\textbf{Connection.} A useful conceptual check on proposed tube guarantees:
state the null being tested, the sampling symmetry, and the coverage event.
A valid test for a future random outcome does not by itself yield a simultaneous
band for the conditional mean reward. The impossibility results motivate restricted
subgroups, explicit structural assumptions, or approximate conditional guarantees.}

\paper{CQR and minimax limits under known covariate shift \cite{cqrlimits}}{
\textbf{Theoretical detail.} The paper analyzes conditional coverage and length
errors relative to oracle equal-tailed intervals, separating quantile-estimation,
calibration, and rank-discretization effects. It assumes known bounded density
ratios in its shift results. In its explicit fixed-score, single-carrier
construction, the minimax calibration scale is
\[
 \mathcal R_{m,p}\asymp
 \sqrt{\frac{1+\chi^2(Q_X\Vert P_X)}{m}}.
\]
Its multi-atom construction adds a complexity factor corresponding to the number
of atoms. This suggests the information scale
$m_{\rm eff}=m/(1+\chi^2(Q_X\Vert P_X))$.

\textbf{Conditions and limits.} These are results for the paper's specified
classes and fixed-score calibration setting. They are not a universal lower bound
for every learned representation, reward estimator, or adaptive PRM procedure.
A marginal guarantee alone does not bound the realized conditional coverage profile.

\textbf{Connection.} Helps quantify why highly concentrated optimized policies
can demand substantially more calibration data. A theoretical project could
specialize these error decompositions to problem-level trajectory scores or
selected-prefix populations, while explicitly checking bounded-ratio and score
assumptions. The divergence-based effective sample size is a diagnostic, not an
independent proof of trajectory coverage.}

\section{Revised literature}
The following literature supplies the established methods, theoretical foundations,
and evaluation resources needed to interpret the recent developments. Entries retain
their mathematical details and limitations; papers already discussed above are not repeated.

\subsection{Closest predecessors and the possible contribution}
Uncertain rewards and pessimistic optimization are established ideas. Four
comparisons are especially important:
\begin{enumerate}
\item AdvPO constructs a shared last-layer confidence set and optimizes
reference-centered pessimistic improvement \cite{advpo}.
\item Quantile and conditional-OT PRMs estimate prefix continuation-success
uncertainty \cite{park,otprm}.
\item Ensembles and reward-distribution DRO yield conservative effective rewards
\cite{coste,rrlhf,unifying}.
\item JOMI addresses prediction coverage after data-dependent selection
\cite{jomi}; correlation-constrained robust RL already supplies structured
reward sets and safe-improvement bounds \cite{correlated,robustcorr}.
\end{enumerate}
Our proposed contribution must therefore concern \emph{process-specific targets,
shared error geometry, coverage after optimization, or an informative failure
characterization}, rather than the generic idea of placing a tube around a reward.

\subsection{Process supervision: targets, data, and representation}

\paper{Lightman et al., Let's Verify Step by Step (2023)}{
\textbf{Model and objective.} Human annotators label steps as positive, negative,
or neutral. The PRM predicts label probabilities conditioned on the problem and
prefix; training is step-label cross-entropy. The preferred solution score in
Appendix F is
\[
 S(\tau)=\prod_{t=1}^{T}p_\theta(\text{positive or neutral}\mid h_t).
\]
The paper compares this with minimum aggregation and different handling of
neutral labels, and evaluates best-of-$N$ selection on mathematical solutions
\cite{lightman}.

\textbf{Conditions and evidence.} PRM800K supplies human process supervision;
active learning concentrates annotation on informative candidates. Annotation
protocols and termination at a first error affect which prefixes receive labels.
The product is a ranking rule: it is a joint-validity probability only if its
factors have the appropriate conditional-event interpretation. It also penalizes
longer trajectories.

\textbf{Connection.} Use this as the intrinsic-validity data anchor and nominal
PRM baseline. Our tube must define its treatment of neutral labels and censored
post-error steps. Strong candidate-selection performance does not establish
coverage or robustness under an RL-trained optimizer.}

\paper{Wang et al., Math-Shepherd (ACL 2024; arXiv v2)}{
\textbf{Model and objective.} A completer samples $K$ continuations from each
prefix. With $C_k$ indicating a correct final answer, the two annotation rules are
\[
 \tilde y_t^{\rm SE}=K^{-1}\sum_k C_k,\qquad
 \tilde y_t^{\rm HE}=\ind\{\max_k C_k=1\}.
\]
The verifier learns these labels and supports reranking and process-supervised
PPO; see the annotation section and Eqs. (3)--(4) \cite{shepherd}.

\textbf{Conditions and evidence.} Under independent Bernoulli continuations,
soft estimation has mean $p_{\pi_c}(h_t)$ and variance $p_{\pi_c}(1-p_{\pi_c})/K$.
Our calculation gives $\E\tilde y_t^{\rm HE}=1-(1-p_{\pi_c}(h_t))^K$.
Thus hard labels depend on both completer capability and rollout count; they
are not unbiased estimates of the single-rollout success probability. A repaired
incorrect step can receive a positive label.

\textbf{Connection.} Separate finite-rollout noise from incorrect target choice.
Increasing $K$ shrinks estimation noise for the soft target but does not make
completion success equal logical correctness. Freeze the completer when
calibrating a success-probability tube.}

\paper{Zhang et al., The Lessons of Developing PRMs (2025; arXiv v1)}{
\textbf{Method.} The study compares Monte Carlo labels, LLM judgments, and human
annotations. Consensus filtering retains solutions when MC and the judge agree
on the first erroneous step. In consistent notation, its selection condition is
\[
 \mathcal D_{\rm keep}=\{\tau:k_{\rm MC}(\tau)=k_{\rm judge}(\tau)\}.
\]
The trained PRMs use step classification; the study also compares soft MC labels
with hard labels $y_t=\ind\{\tilde p_t>0\}$ and evaluates both best-of-$N$ and
ProcessBench \cite{lessons}.

\textbf{Conditions and evidence.} In the reported eight-completion experiments,
hard labels perform particularly well after filtering. Agreement is a quality
filter, not proof that both annotators are correct. Conditioning on agreement
also changes the data distribution and can discard difficult cases. Final-answer
ranking and first-error identification need not improve together.

\textbf{Connection.} Keep corrected supervision identical across uncertainty
methods. Audit the rejected and retained strata: a tube calibrated only on
consensus examples may fail on precisely the disputed prefixes that an optimizer
finds. This paper supplies practical Qwen PRM comparators, not a confidence-set
guarantee.}

\paper{Ding et al., SCAN (2025; arXiv v1)}{
\textbf{Method and formula.} SCAN estimates problem-level self-confidence
$SC_\pi(x)$ as the fraction of correct sampled answers, uses high-confidence
positive examples, and concentrates stepwise MC annotation on negatives.
Near the estimated first error $k_e$, its soft-label correction has the form
\[
 \hat y_t=\min\{c_t/SC_\pi(x),1\};
 \qquad \hat y_t=\ind\{c_t>0\}\ \text{away from that region},
\]
where $c_t$ is continuation success. A tolerance distance controls the affected
steps; see Section 3.2, Eqs. (5)--(7) \cite{scan}.

\textbf{Conditions and evidence.} Confidence-based filtering reduces annotation
cost, while soft labels near a boundary accommodate error-location noise.
Division by self-confidence is a heuristic capability correction, not an
identified importance ratio or an unbiased estimator of intrinsic validity.
Its denominator must be nonzero; filtering is part of the protocol.

\textbf{Connection.} A useful comparator for improving the reward center and
annotation efficiency. Stratify tube errors by problem confidence and distance
to the estimated error. Check difficult low-confidence problems independently,
since training selection can hide their uncertainty.}

\paper{Pala et al., Error Typing for Smarter Rewards / PathFinder-PRM (2025; arXiv v1)}{
\textbf{Model and objective.} PathFinder decouples error detection from reward
estimation using two passes:
\[
 (M_t,C_t)=f_\theta(h_t),\qquad
 R_t=g_\theta(h_t,M_t,C_t).
\]
The first predicts mathematical and consistency signals; the second conditions
its reward on those signals. Training uses cross-entropy on categorical label
tokens for mathematical validity, consistency, and step correctness
\cite{pathfinder}.

\textbf{Conditions and evidence.} Section 3.2 calls the dataset approximately
400K trajectory samples, assembled from PRM800K and Mistral-PRM data; the precise
serialized sample unit matters for annotation-cost comparisons. The paper maps
neutral PRM800K labels into a category interpreted as suboptimal, which is its
labeling convention rather than a universal definition of neutral validity.
Both passes may share representation and judgment errors.

\textbf{Connection.} Typed residuals can define distinct tube directions or
coverage strata. Test whether any gain comes from the uncertainty geometry or
simply richer labels. Account for the second forward pass when matching compute;
error-type predictions do not themselves certify reward bounds.}

\subsection{Credit assignment and adversarial failure discovery}

\paper{Cheng et al., Stop Summation / PURE (NeurIPS 2025; arXiv v1)}{
\textbf{Objective and implementation.} PURE targets
$\max_\pi\E_\pi[\min_t r_t]$, rather than a sum. Its transformed step reward is
\[
 \omega_t=\frac{\exp(-r_t/\vartheta)}{\sum_j\exp(-r_j/\vartheta)},\qquad
 r_t^{\rm tr}=\omega_t r_t.
\]
For a unique worst step $w$, $\vartheta\downarrow0$ concentrates reward at $w$.
With $\gamma=1$, return-to-go is $r_w$ before/at $w$ and zero afterwards.
Section 3.2 adds a leave-one-out baseline and optionally verifiable outcome
reward \cite{pure}.

\textbf{Conditions and evidence.} The exact minimum correspondence is a
zero-temperature limit; finite temperature is a smooth reweighting. Ties share
weight across worst steps. Minimum aggregation limits accumulation of favorable
errors but can reward vacuous outputs with no low-scoring mathematical claims.
It also treats a weak step differently from a recoverable reasoning error.

\textbf{Connection.} Required aggregation/credit-assignment baseline. Hold
$G$ fixed before comparing uncertainty estimators. Our extension could apply
minimum aggregation to lower bounds, but any improvement must exceed simply
switching from sum to minimum or adding an outcome check.}

\paper{Zheng et al., CoLD (arXiv v2, May 2026)}{
\textbf{Model and objective.} A noisy-feature bias estimator gives
$q_\phi(h)=\sigmoid(b_\phi(h+\mathcal N))$. The inference score in Eq. (3) is
\[
 r_{\rm deb}(h)=r_\theta(h)q_\phi(h)-c q_\phi(h)-\alpha\ell(h),
\]
where $\ell$ measures step length. Joint training uses correctness cross-entropy
and complementary correlation terms
$\lambda_r\operatorname{Corr}(r_\theta,\ell)^2$ and
$-\lambda_b\operatorname{Corr}(b_\phi,\ell)^2$ \cite{cold}.

\textbf{Conditions and evidence.} The estimator is intended to capture length
and surface style; Gaussian feature noise approximates the counterfactual
removal of semantics. That approximation does not identify a causal length
effect without further assumptions. The paper's limitation discussion recognizes
this issue. A deterministic length penalty also changes preferences between
long and short correct solutions.

\textbf{Connection.} Length is a concrete common error direction for a structured
tube. Compare against this direct bias correction and measure residual errors
under meaning-preserving verbosity changes. A learned bias parameter's uncertainty
would need independent estimation; neither the correlation penalty nor injected
noise supplies confidence coverage.}

\paper{Wallace et al., Universal Adversarial Triggers (EMNLP 2019)}{
\textbf{Attack construction.} A short shared token sequence is optimized to
cause a target behavior across many inputs. The objective can be written
$\min_u\E_{x\sim\mathcal D}\mathcal L(f(u\mathbin{\oplus}x),y_{\rm target})$.
A first-order token replacement ranks vocabulary embeddings using
\[
 v\in\arg\min_{v'\in\mathcal V}
 (e_{v'}-e_{u_j})^\top\nabla_{e_{u_j}}\mathcal L,
\]
then evaluates candidate replacements and iterates \cite{triggers}.

\textbf{Conditions and evidence.} This is a gradient-based discrete search
approximation, not a proof that the chosen trigger globally optimizes the loss.
The original demonstrations concern classification, reading comprehension, and
language generation; transfer to a particular PRM is an empirical question.
A universal trigger differs from optimizing a separate attack for each problem.

\textbf{Connection.} Test shared prefixes, suffixes, and reward-delimiter
perturbations in addition to semantic mistakes. Shared attacks can reveal a
common model bias that ensemble spread misses. Keep the permitted token
locations explicit; formatting attacks may exploit the scorer's parser rather
than uncertainty about the mathematical reasoning itself.}

\subsection{Reward hacking theory and optimization pressure}

\paper{Skalse et al., Defining and Characterizing Reward Hacking (arXiv v2, 2025)}{
\textbf{Definition.} Proxy $\hat r$ and target $r^*$ are hackable over policy
class $\Pi$ if there exist $\pi,\pi'\in\Pi$ with
\[
 J_{\hat r}(\pi')>J_{\hat r}(\pi),\qquad
 J_{r^*}(\pi')<J_{r^*}(\pi).
\]
Additive returns are linear in discounted visit counts:
$J_r(\pi)=\langle r,\mathcal F^\pi\rangle$ \cite{skalse}.

\textbf{Theoretical result.} The revised full text's Theorem 1 states that
nontrivial unhackable reward pairs are equivalent on a stationary-policy set
containing an open subset; Corollary 1 covers all stationary policies. Equivalent
means the same policy ordering, not identical numerical rewards. Finite policy
sets permit nontrivial nonequivalent unhackable pairs. The abstract's compressed
wording should not replace these precise statements.

\textbf{Connection.} Aim for restricted-domain coverage or safe improvement
relative to a reference, rather than universal nontrivial proxy alignment.
A finite candidate pool has a different guarantee scope from all policies.
Our tube can reduce specific reversals without eliminating every possible
reversal or implying that every nominal RL run must hack.}

\paper{Gao et al., Scaling Laws for Reward Model Overoptimization (2022)}{
\textbf{Empirical formulas.} With $d=\sqrt{\KL(\pi\|\pi_0)}$, gold-model
reward improvement is fitted by
\[
 R_{\rm BoN}(d)=d(\alpha_{\rm B}-\beta_{\rm B}d),\qquad
 R_{\rm RL}(d)=d(\alpha_{\rm R}-\beta_{\rm R}\log d).
\]
Coefficients vary with reward-model size and data. For continuous, untied
best-of-$N$ selection, the paper uses
$D_{\mathrm{KL}}^{\rm BoN}=\log N-(N-1)/N$ \cite{gao}.

\textbf{Conditions and evidence.} These are fitted scaling relationships,
not universal laws. The RL expression is not a good local description at the
origin. A larger learned gold model supplies synthetic preferences and evaluation;
it is still a proxy for human intent. The analytic BoN KL formula assumes the
particular independent-sampling/ranking model and needs care with ties or
restricted discrete supports.

\textbf{Connection.} Plot target quality against $N$, training steps, and
policy divergence, alongside proxy reward. Compare the point at which quality
peaks and declines. Mathematical answer checks and independent process audits
can replace a learned gold reward where available; one favorable checkpoint
cannot establish mitigation.}

\paper{Laidlaw et al., Correlated Proxies (ICLR 2025; arXiv v3)}{
\textbf{Definition and theorem.} Let
$\operatorname{Corr}_{d_{\pi_0}}(\hat r,r^*)=\rho>0$ and let
$\sigma_{\hat r},\sigma_{r^*}$ be reference-distribution standard deviations.
For $d_\pi\ll d_{\pi_0}$, Theorem 5.1 gives
\[
 \frac{\Delta J_{r^*}}{\sigma_{r^*}}
 \geq\frac{1}{\rho}\left[
 \frac{\Delta J_{\hat r}}{\sigma_{\hat r}}
 -\sqrt{(1-\rho^2)\chi^2(d_\pi\|d_{\pi_0})}\right].
\]
This motivates occupancy-regularized policy optimization (ORPO)
\cite{correlated}.

\textbf{Conditions.} Correlation is under the reference occupancy, not a
pointwise error guarantee. Nonzero variances, support overlap, and normalized
occupancy conventions matter. A positive lower bound certifies improvement;
there need not exist a policy achieving such a bound. Action-level KL and
occupancy divergence are different objects.

\textbf{Connection.} Relate tube residual error to the change in prefix
occupancy, and test whether policy movement explains apparent mitigation.
Estimating a correlation or density ratio in enormous text state spaces is
itself a statistical problem. This result supplies a comparison for any proposed
reference-centered PRM safe-improvement theorem.}

\subsection{Algorithmic context for the experiments}

\paper{Sutton and Barto, Reinforcement Learning: An Introduction (2018)}{
\textbf{Core equations.} For a fixed policy in a discounted MDP, the Bellman
expectation equation is
\[
 V^\pi(s)=\sum_a\pi(a\mid s)\sum_{s',r}p(s',r\mid s,a)
 [r+\gamma V^\pi(s')].
\]
Temporal-difference learning uses
$\delta_t=r_t+\gamma V(s_{t+1})-V(s_t)$ and
$V(s_t)\leftarrow V(s_t)+\alpha\delta_t$; the policy-gradient theorem underlies
actor--critic learning \cite{sutton}.

\textbf{Conditions.} The state must contain enough history for the Markov
property. Infinite-horizon discounted values require appropriate boundedness
and $\gamma<1$; finite-horizon values also depend on time. Standard Bellman
recursions assume additive rewards and a fixed model. Approximate values and
bootstrapping introduce estimation error separate from PRM error.

\textbf{Connection.} Use Part II to translate these objects into text prefixes,
tokens/steps, and terminal correctness. Our tube changes the reward objective,
while a critic estimates return under that objective. Confusing critic variance
with reward uncertainty would obscure the mechanism. This is a textbook
background reference, not a new reward-hacking result.}

\paper{Schulman et al., Proximal Policy Optimization (2017)}{
\textbf{Objective.} With importance ratio
$q_t(\theta)=\pi_\theta(a_t\mid s_t)/\pi_{\rm old}(a_t\mid s_t)$,
PPO's clipped surrogate is
\[
 L^{\rm clip}(\theta)=\hat\E_t\min\{q_t\hat A_t,
 \clip(q_t,1-\epsilon,1+\epsilon)\hat A_t\}.
\]
The full training objective can add value-fitting and entropy terms; the paper
also studies adaptive KL penalties \cite{ppo}.

\textbf{Theoretical distinction.} The minimum is a lower surrogate relative
to the unclipped sample objective. It is not a confidence bound on true reward
and does not enforce a global trust region for every state. Multiple epochs
reuse on-policy samples while controlling update incentives; a bad reward
still gives a bad advantage estimate.

\textbf{Connection.} Keep clipping, critic fitting, KL coefficients, and
optimization epochs fixed when comparing nominal and robust PRM rewards.
Recompute the intended robust return/advantage consistently. PPO's use of
``pessimism'' concerns an optimization surrogate and must not be conflated with
pessimism over an uncertain reward function.}

\paper{Shao et al., DeepSeekMath / GRPO (2024; arXiv v3)}{
\textbf{Objective and credit assignment.} GRPO eliminates a learned value
baseline by sampling a group and normalizing its rewards. Outcome advantages are
\[
 \hat A_i=\frac{R_i-\bar R}{s_R};\qquad
 \hat A_{i,t}^{\rm process}=\sum_{j:\operatorname{end}(j)\geq t}
 \frac{r_{ij}-\bar r}{s_r}.
\]
These feed a PPO-style clipped objective with a separate policy KL term.
The process formula normalizes over group step rewards; see Section 4.1
\cite{deepseekmath}.

\textbf{Conditions and evidence.} Groups with identical outcomes give zero
centered advantages and require a zero-standard-deviation convention. Summing
future normalized process rewards changes length incentives and credit
assignment. Group normalization can also amplify small differences. The paper
studies iterative reward-model updates; the target distribution can consequently
change during training.

\textbf{Connection.} A natural RL testbed, but a lower PRM score is not by
itself a lower bound on the normalized advantage. Audit how our tube's scores
pass through normalization and aggregation. Distinguish reward-model uncertainty,
Monte Carlo group variability, and policy-gradient estimation noise.}

\subsection{Uncertainty and pessimism: established predecessors}

\paper{Coste et al., Reward Model Ensembles (arXiv v2, 2024)}{
\textbf{Aggregation rules.} For $K$ reward models, define
$\mu=K^{-1}\sum_k r_k$ and $v=K^{-1}\sum_k(r_k-\mu)^2$. The compared scores are
\[
 r_{\rm mean}=\mu,\qquad r_{\rm WCO}=\min_k r_k,
 \qquad r_{\rm UWO}=\mu-\lambda v.
\]
Worst-case optimization (WCO) and uncertainty-weighted optimization (UWO)
are tested with best-of-$N$ and PPO; see Section 3, Eqs. (4)--(5)
\cite{coste}.

\textbf{Evidence and assumptions.} Experiments use a larger learned gold
reward, including synthetic preference-label flips. Member diversity comes
from initialization and training randomness on shared data. Minimum-member
scoring is below truth at a point if at least one member is not overestimating
there; the method does not guarantee that condition. Its pointwise minimum
also need not correspond to one shared reward model for all trajectories.

\textbf{Connection.} Include all three rules with matched ensemble cost.
Distinguish variance from standard-deviation penalties and align reward scales.
Test common errors rather than treating member spread as calibrated confidence.
Our coupled function set should be compared with both per-sample WCO and a
single worst ensemble member chosen at policy level.}

\paper{Zhang et al., AdvPO (2024; arXiv v2)}{
\textbf{Confidence set.} With frozen final features $e(x,y)$,
$r_\phi=\phi^\top e$, construct
\[
 M_D=\lambda I+\sum_{i,y\in\{y_i^+,y_i^-\}}e(x_i,y)e(x_i,y)^\top,
 \qquad \mathcal C=\{\phi:\norm{\phi-\hat\phi}_{M_D}\leq b\}.
\]
For $g=\E_\pi e-\E_{\rm ref}e$, reference-centered robust improvement becomes
\[
 \inf_{\phi\in\mathcal C}\phi^\top g
 =\hat\phi^\top g-b\sqrt{g^\top M_D^{-1}g}.
\]
A policy KL penalty completes the objective; Section 4, Eqs. (5)--(7),
Theorem 4.1 \cite{advpo}.

\textbf{Theory and limits.} Theorem 3.1's reward bound includes an approximation
term and assumes infinite-width/positive-definite NTK structure. Experimental
radii are tuned rather than independently coverage-calibrated. One common
parameter adversary is less conservative than separate worst-case choices for
each sample, by convexity of the norm; this distinction is already central here.

\textbf{Connection.} The closest predecessor to Approach B and to
reference-centered safe improvement. A publishable extension must add
process-specific construction, measured joint coverage, or misspecification
control, rather than merely repeat the ellipsoid's dual-norm formula.}

\paper{Yan et al., Reward-Robust RLHF (2024; arXiv v1)}{
\textbf{Objective and uncertainty model.} For a finite reward family $\U$,
\[
 J_\lambda(\pi)=\lambda J_{\hat r}(\pi)
 +(1-\lambda)\min_{r\in\U}J_r(\pi).
\]
The Bayesian Reward Model Ensemble shares a backbone and has Gaussian heads
$r_k=\mu_k+\sigma_k\varepsilon$, with $\varepsilon\sim\mathcal N(0,1)$.
The nominal head is selected using low uncertainty; see Sections 5--5.1
\cite{rrlhf}.

\textbf{Conditions and evidence.} The nominal/robust tradeoff reduces excessive
conservatism but weakens a purely worst-case interpretation. Shared features
can produce common bias, and Gaussian sampling is a modeling assumption,
not proof that the true reward lies in the family. The paper's constant-reward
argument requires action-independent remaining return: variable termination or
length-dependent cumulative constants need separate treatment. Policy KL can
still move the actor even when reward advantages vanish.

\textbf{Connection.} Compare robust tubes with a nominal--robust interpolation,
but report the tradeoff parameter and coverage independently. Preserve the
minimum-outside-expectation objective when implementing a shared-function
adversary; swapping that order changes the method and its conservatism.}

\paper{Park et al., Know What You Don't Know (2025; arXiv v1)}{
\textbf{Target and calibration.} Fixed-completer rollouts produce
$\tilde p(h)=K^{-1}\sum_k C_k$. Quantile heads minimize
\[
 \mathcal L_{\rm QR}=\frac{1}{|\mathcal B|}\sum_{\beta\in\mathcal B}
 \rho_\beta\bigl(\tilde p(h)-q_\beta(h)\bigr).
\]
Instance-adaptive scaling uses a conservative predicted success value $p_L$:
\[
 N_{\rm IAS}=\min\!\left\{N_{\max},
 \left\lceil\frac{\log(1-C)}{\log(1-p_L)}\right\rceil\right\}.
\]
See Sections 4.1--4.2 \cite{park}.

\textbf{Conditions and theory.} The budget identity uses independent trials:
$1-(1-p)^N\geq C$. The paper's Theorem 1 considers an uncapped budget and a
held-out calibration adjustment. Fitted quantile regression alone is not exact
finite-sample conformal coverage. A cap can defeat the target confidence.
The existence of a correct candidate does not guarantee that the reranker
chooses it; $p_L=0$ requires a maximum-budget convention.

\textbf{Connection.} A required calibrated-PRM baseline. Its uncertainty target
is continuation success, not intrinsic validity. Our extension should test
selected-candidate and simultaneous-prefix coverage, rather than assuming that
marginal calibration survives search or policy updates.}

\subsection{Statistical foundations for constructing and calibrating a tube}

\paper{Deep ensembles \cite{deepensemble}}{
\textbf{Model and equations.} For Gaussian regression, member $m$ minimizes
$\frac12\log\sigma_m^2(x)+\frac{(y-\mu_m(x))^2}{2\sigma_m^2(x)}$.
The predictive mixture has moments
\[
 \bar\mu=\frac1M\sum_m\mu_m,\qquad
 v=\frac1M\sum_m\sigma_m^2+\frac1M\sum_m(\mu_m-\bar\mu)^2.
\]
The second term measures disagreement between fitted means; the first represents
within-member predictive variation. These are mixture moments, not a claim that
the mixture itself is Gaussian. The paper combines proper scoring rules with
independently initialized networks; classification predictions average probabilities.

\textbf{Conditions and evidence.} Predictive calibration and robustness are
empirical findings. Independent initialization does not establish independence
of errors on new inputs, especially with shared training data or representations.
An ensemble standard deviation is neither a finite-sample confidence radius nor
a guarantee of coverage under optimization.

\textbf{Connection.} A practical starting estimator for the tube center and
heterogeneous width. Retaining each member's entire vector of prefix predictions
also supplies coherent trajectory scenarios. Compare between-member uncertainty
against total predictive variance: penalizing irreducible outcome randomness can
reject valid but difficult reasoning. Independent calibration and attack tests
remain necessary.}

\paper{Self-normalized confidence sets for linear bandits \cite{abbasi}}{
\textbf{Theorem and equations.} With $Y_t=X_t^\top\theta_*+\eta_t$,
predictable features, conditionally $\sigma$-sub-Gaussian noise, and
$\norm{\theta_*}_2\le S$, define
\[
 V_n=\lambda I+\sum_{t\le n}X_tX_t^\top,\quad
 \widehat\theta_n=V_n^{-1}\sum_{t\le n}X_tY_t.
\]
Theorem 2 gives, simultaneously over $n$, with probability at least $1-\delta$,
\[
 \norm{\widehat\theta_n-\theta_*}_{V_n}
 \le \sigma\sqrt{2\log\frac{\det(V_n)^{1/2}}
 {\det(\lambda I)^{1/2}\delta}}+\sqrt\lambda S=:b_n.
\]
Thus prediction error is bounded by $b_n\norm{\phi(h)}_{V_n^{-1}}$ on that event.

\textbf{Conditions and scope.} Adaptively chosen features are allowed, but the
linear model, predictable design, noise condition, and parameter bound are
substantive assumptions. A neural feature map learned on the same data, a
misspecified linear head, or logistic preference likelihood requires additional
analysis. OFUL uses the set for optimism, rather than pessimistic reward selection.

\textbf{Connection.} This is a rigorous template for a shared-parameter tube,
with larger widths in poorly observed directions. AdvPO already adapts related
geometry to reward optimization. A project contribution must address PRM-specific
dependence, misspecification, or coverage, rather than rediscover this ellipsoid.}

\paper{Conformalized quantile regression (CQR) \cite{cqr}}{
\textbf{Construction and guarantee.} Fit lower and upper quantiles on a training
split. On $m$ held-out examples compute
\[
 S_i=\max\{\widehat q_{\alpha/2}(X_i)-Y_i,
 Y_i-\widehat q_{1-\alpha/2}(X_i)\},\quad
 k=\lceil(m+1)(1-\alpha)\rceil.
\]
Let $q$ be the $k$th calibration order statistic, taking $q=+\infty$ if
$k>m$. The interval is
\[
 C(x)=[\widehat q_{\alpha/2}(x)-q,
        \widehat q_{1-\alpha/2}(x)+q].
\]
Exchangeability of calibration and test examples yields marginal coverage
$\Prob\{Y_{\rm new}\in C(X_{\rm new})\}\ge1-\alpha$;
the upper bound is $1-\alpha+1/(m+1)$ under appropriate absence of ties.

\textbf{Conditions and limits.} The quantile model need not be correct for
marginal coverage, but affects efficiency. The guarantee does not assert
coverage conditional on every prefix, simultaneously over a trajectory, or after
reward-based selection. Binary targets can produce uninformative prediction sets.

\textbf{Connection.} CQR calibrates a future random label or return, not directly
the unknown mean reward function. To obtain a trajectory tube, define a suitable
trajectory-level score and calibration unit; do not simply reinterpret stepwise
prediction intervals as simultaneous confidence bands.}

\paper{JOMI: selection-conditional conformal prediction \cite{jomi}}{
\textbf{Construction and guarantee.} For a selected test index $j$, JOMI forms
candidate-label-dependent reference sets $\mathcal R_j(y)$ by swapping calibration
and test examples and retaining swaps that preserve the relevant selection event.
Schematically, its prediction set is
\[
 C_j=\left\{y:V(X_j,y)\le
 Q_{1-\alpha}\bigl(\{V_i:i\in\mathcal R_j(y)\}\cup\{+\infty\}\bigr)\right\}.
\]
Under the paper's permutation and exchangeability conditions, coverage holds
conditional on $j$ being selected and a specified class of selection outcomes.
Randomization can handle ties and sharpen exactness.

\textbf{Conditions and limits.} The selection rule must satisfy the required
invariance, and the reference set must actually implement that rule. The relevant
calibration/test exchangeability is conditional on other test covariates.
Candidates generated for one shared problem are dependent; treating every prefix
or candidate as an independent calibration example does not establish this premise.

\textbf{Connection.} Especially relevant to best-of-$N$, where ordinary marginal
coverage can fail at the selected maximum. Problem-level calibration or a justified
conditional sampling model is needed. This covers the selected prediction target;
it does not automatically give joint trajectory coverage or validity throughout
an adaptive RL training loop.}

\paper{Weighted conformal prediction under covariate shift \cite{weighted}}{
\textbf{Construction and guarantee.} Suppose $Q_X\ll P_X$, the density ratio
$w(x)=dQ_X/dP_X$ is known, and $P_{Y\mid X}=Q_{Y\mid X}$. Split conformal replaces
the ordinary score quantile by the weighted distribution
\[
 \sum_{i=1}^m\frac{w(X_i)}{\sum_{k=1}^m w(X_k)+w(x)}\delta_{S_i}
 +\frac{w(x)}{\sum_{k=1}^m w(X_k)+w(x)}\delta_{+\infty}.
\]
Its $(1-\alpha)$ quantile yields marginal coverage for a test point drawn from
$Q$, while calibration examples come from $P$.

\textbf{Conditions and limits.} The infinite test-point mass is essential to
the finite-sample correction. Poor overlap can force an infinite quantile.
Estimated or clipped weights do not inherit the exact theorem without accounting
for their error. The theorem assumes stable conditional labels, not arbitrary
joint distribution shift.

\textbf{Connection.} This offers a route from reference-policy calibration to a
new policy's prefix distribution. If a PRM target is continuation success, changing
the rollout policy can change $Y\mid X$, violating covariate shift. Policy selection
using calibration labels also needs separate treatment. Measure overlap and weight
concentration before claiming a usable tube under optimization.}

\paper{Conformal prediction beyond exchangeability \cite{beyond}}{
\textbf{Theoretical detail.} The paper quantifies degradation using swapped-data
laws rather than assuming exact exchangeability. A representative coverage-gap
bound for prespecified weights $w_i\in[0,1]$ has the form
\[
 (1-\alpha)-\Prob\{Y_{n+1}\in C(X_{n+1})\}
 \le\frac{\sum_{i=1}^n w_i\,
 \operatorname{TV}(\mathcal L(Z),\mathcal L(Z^{(i)}))}
 {1+\sum_{i=1}^n w_i},
\]
where $Z^{(i)}$ swaps observation $i$ with the test observation. The paper also
handles nonsymmetric fitting procedures using its randomization construction.

\textbf{Conditions and limits.} The right side must be small to obtain a useful
statement. Recency weighting can reduce the influence of old observations but
also reduces calibration information. Data-dependent weights require the relevant
conditioning or additional justification. Bounding swap distances for adaptively
generated reasoning trajectories is a separate research problem.

\textbf{Connection.} Useful for analyzing gradual policy drift and periodically
recalibrated PRMs. It supplies a degradation statement, not a guarantee under
arbitrary reward-driven distribution shift. An experiment should report how coverage
changes with policy divergence and calibration age, rather than silently assuming
that a reference-policy conformal guarantee survives RL.}

\paper{Conformal risk control (CRC) \cite{crc}}{
\textbf{Construction and theorem.} Let $L_i(\lambda)$ be exchangeable loss
functions, bounded above by $B$, nonincreasing and right-continuous in $\lambda$.
With $\widehat R_n(\lambda)=n^{-1}\sum_iL_i(\lambda)$, choose
\[
 \widehat\lambda=\inf\left\{\lambda:
 \frac{n}{n+1}\widehat R_n(\lambda)+\frac{B}{n+1}\le\alpha\right\}.
\]
With an admissible conservative endpoint and the paper's conditions,
$\E[L_{n+1}(\widehat\lambda)]\le\alpha$.

\textbf{Conditions and limits.} This is an expected-risk guarantee, not a
high-probability bound on every deployed instance or risk conditional on selection.
Monotonicity concerns the actual loss function. If changing $\lambda$ changes which
trajectory is selected, the resulting error rate need not be monotone. Empty-set
or abstention behavior must be included in the definition.

\textbf{Connection.} A suitable PRM application is thresholding fixed scores:
\[
 L_i(\lambda)=\ind\{\text{invalid step}_i\}\ind\{s_i\ge\lambda\}.
\]
This loss is nonincreasing. This controls erroneous acceptance frequency, which differs from
the fraction of accepted steps that are invalid. For a pessimism coefficient used
inside best-of-$N$ or RL, first prove the required monotonicity or use another
calibration method.}

\paper{CROMS: model selection for conformalized robust optimization \cite{croms}}{
\textbf{Decision formulation.} A conformal set induces the robust decision
\[
 z_\lambda(x)\in\arg\min_z\max_{y\in C_\lambda(x)}\varphi(y,z),\qquad
 \lambda_*\in\arg\min_\lambda\E[\varphi(Y,z_\lambda(X))].
\]
Thus model selection targets realized decision loss, rather than interval width
alone. The paper distinguishes efficient E-CROMS, with approximate/asymptotic
validity and optimality results, from a symmetric augmented-data F-CROMS construction
with finite-sample coverage and asymptotic optimality. Its localized variant
requires additional nonparametric assumptions.

\textbf{Conditions and limits.} Ordinary split-conformal coverage cannot simply
be retained after arbitrarily choosing the model on the same calibration data.
The guarantees depend on the specified selection construction and regularity
conditions. Optimizing a narrow interval can select a model with worse robust
decisions; conversely, conservative decisions can lose substantial nominal utility.

\textbf{Connection.} Particularly relevant when choosing among tube geometries,
uncertainty estimators, or pessimism levels. Calibrate and select the entire
uncertainty-to-decision pipeline, and evaluate the resulting reasoning success and
abstention cost. Decision-aware calibration is already established here; novelty
would require a PRM-specific selection/dependence analysis or demonstrably improved
robust reasoning behavior.}

\subsection{Robust optimization and reward shaping}

\paper{Distributionally robust optimization (Duchi--Namkoong) \cite{duchi}}{
\textbf{Objective and duality.} For loss $\ell_\theta$, optimize
\[
 \min_\theta\sup_{Q\ll P:\,D_f(Q\Vert P)\le\rho}\E_Q[\ell_\theta].
\]
Under suitable duality conditions, its inner value can be expressed as
\[
 \inf_{\lambda\ge0,\eta}\left\{
 \lambda\rho+\eta+\lambda\E_P
 \left[f^*\left(\frac{\ell_\theta-\eta}{\lambda}\right)\right]\right\},
\]
with boundary cases interpreted by limits. The paper develops guarantees for
specified divergence families, including Cressie--Read divergences, under loss
and statistical regularity assumptions.

\textbf{Conditions and limits.} Robustness to a subpopulation follows only when
its conditional distribution lies inside the ambiguity set. A ball around an
empirical distribution reweights observed support; it cannot create an unseen
failure type. Divergence choice, radius, and tail assumptions determine how
conservative and statistically estimable the objective is.

\textbf{Connection.} Distinguish uncertainty about the distribution of problems
or trajectories from uncertainty about the reward function itself. They can be
combined, but their adversaries and coverage events differ. A PRM project could
use DRO for rare reasoning patterns while a reward-function tube handles evaluator
error; identifying which uncertainty actually causes hacking is essential.}

\paper{Robust dynamic programming (Iyengar) \cite{iyengar}}{
\textbf{Structural principle.} Rectangular transition ambiguity permits local
adversarial choices compatible with dynamic programming. A standard state-action
rectangular discounted recursion can be represented as
\[
 (Tv)(s)=\max_a\inf_{p\in\mathcal P_{s,a}}
 \left\{r(s,a)+\gamma\sum_{s'}p(s')v(s')\right\}.
\]
For bounded rewards, $\gamma<1$, and suitable nonempty uncertainty sets, the local
operator preserves the contraction structure needed for value iteration. The
paper's central contribution is the role of rectangularity in robust dynamic
programming and associated game interpretations.

\textbf{Conditions and limits.} Rectangularity means the uncertainty set allows
the relevant combinations of local choices. A single globally shared uncertain
parameter typically couples states and violates that structure. Substituting
independent local worst cases changes the adversary and can produce excessive
pessimism. This displayed recursion illustrates the standard specialization;
it is not a claim that every reward tube meets the paper's assumptions.

\textbf{Connection.} A trajectory tube must specify whether errors can be chosen
independently at every prefix or must arise from one coherent reward function.
That choice determines both conservatism and computational tractability. Robust
Bellman updates require an explicit rectangular formulation or a justified state
augmentation, not merely an uncertainty interval attached to each step.}

\paper{Potential-based reward shaping \cite{shaping}}{
\textbf{Theorem and telescoping identity.} The shaping reward
\[
 F(s,a,s')=\gamma\Phi(s')-\Phi(s)
\]
changes a finite-horizon discounted return by
\[
 \sum_{t=0}^{T-1}\gamma^t(r_t+F_t)
 =\sum_{t=0}^{T-1}\gamma^t r_t-\Phi(s_0)+\gamma^T\Phi(s_T).
\]
With a fixed initial state and appropriate terminal boundary, or a bounded
potential in an infinite discounted horizon, this preserves policy ordering.
The necessity result concerns transformations that preserve optimal policies
across the paper's broad class of environments.

\textbf{Conditions and limits.} A poor but fixed potential can still preserve
ordering when used consistently; it need not accelerate learning. Terminal
potentials and discount conventions matter. Arbitrary positive process rewards,
stepwise probability products, or min-form returns are not automatically
potential-based shaping.

\textbf{Connection.} If a PRM estimates a reasoning potential, uncertainty should
be placed on a shared potential function and propagated through differences.
Independently minimizing each difference can break telescoping by choosing
incompatible endpoint values. This provides a precise way to distinguish
uncertainty in a policy-invariant shaping signal from uncertainty in a new
objective that can alter the optimal reasoning strategy.}

\subsection{Benchmarks and what their labels can establish}

\paper{ProcessBench \cite{processbench}}{
\textbf{Task and metric.} The benchmark contains approximately 3,400 reasoning
traces, with the earliest erroneous step or a no-error label. Its reported
harmonic metric combines exact first-error localization accuracy on erroneous
traces and no-error accuracy on clean traces:
\[
 F=\frac{2A_{\rm err}A_{\rm clean}}{A_{\rm err}+A_{\rm clean}}.
\]
This is a harmonic mean of two accuracies, rather than ordinary stepwise binary
classification F1. PRM evaluation thresholds are chosen on a designated subset
in the original protocol.

\textbf{Evidence and limits.} The benchmark tests mathematical error detection
across difficulty levels. Correct final answers can coexist with invalid intermediate
reasoning. Earliest-error annotation does not supply verified labels for every
later step, so it cannot directly establish simultaneous all-step coverage.
A static localization score also does not measure robustness to an optimizer
searching for high-scoring wrong solutions.

\textbf{Connection.} Use it to test error localization and clean-trace rejection,
with thresholds and uncertainty calibration fixed independently of evaluation.
Report false acceptance and abstention alongside the benchmark metric. For a
full trajectory-tube coverage claim, obtain all-step labels or define a coverage
event restricted to the observed first-error target.}

\paper{PRMBench \cite{prmbench}}{
\textbf{Task and metric.} PRMBench reports 6,216 instances and 83,456 step labels,
organized around simplicity, soundness, and sensitivity, with nine finer dimensions.
Its PRM-score combines negative-class and overall F1:
\[
 \operatorname{PRM\!\text{-}\!score}
 =w_1F_{1,\rm neg}+w_2F_{1,\rm overall}.
\]
The benchmark probes redundant or circular reasoning, correctness failures, and
sensitivity to changes in reasoning. The weighting and threshold protocol should
be recorded when reproducing results.

\textbf{Evidence and limits.} Step-level discrimination reveals weaknesses
hidden by aggregate final-answer accuracy. However, redundancy is a utility or
efficiency judgment, not necessarily mathematical invalidity. Aggregating these
dimensions into a single binary target can obscure what a PRM's reward represents.
Static test cases do not certify calibration under adversarial optimization.

\textbf{Connection.} Separate correctness, informativeness, and sensitivity in
our labels and reported results. A tube designed for correctness should not be
judged solely by its agreement with a simplicity preference. Paired transformed
traces can test whether uncertainty responds to semantic errors while remaining
stable under harmless rewrites; that invariance diagnostic is a proposed project
evaluation, rather than a coverage theorem supplied by PRMBench.}

\paper{Socratic-PRMBench \cite{socratic}}{
\textbf{Task and metric.} The benchmark has 2,995 samples spanning transformation,
decomposition, regathering, deduction, verification, and integration, with twenty
error types. Its aggregate score is explicitly
\[
 \operatorname{Score}=\tfrac12F_{1,\rm neg}+\tfrac12F_{1,\rm overall}.
\]
The construction uses structured reasoning patterns, synthetic modifications,
and filtering, with a sampled human assessment of annotation agreement.

\textbf{Evidence and limits.} Pattern-level results can reveal that a reward
model recognizes deduction errors but misses verification or integration failures.
Sampled annotation agreement does not certify every synthetic label. Aggregate F1
can conceal a high false-acceptance rate in a rare pattern and says nothing directly
about coverage on optimizer-selected traces.

\textbf{Connection.} A useful stress test for heterogeneous tube widths and
structured uncertainty. Our proposed diagnostic is empirical coverage and false
acceptance within each prespecified pattern, together with interval width and
abstention. Such subgroup results require enough independent problems per group;
they do not imply unrestricted conditional coverage for arbitrary prefixes.
Pair this static benchmark with best-of-$N$ and RL stress tests so that uncertainty
quality is evaluated where reward hacking actually appears.}

\subsection{Synthesis and research positioning}
\begin{longtable}{p{0.20\linewidth}p{0.35\linewidth}p{0.36\linewidth}}
\toprule
Family & Established capability & What our study must resolve \\
\midrule
\endhead
Process supervision & Intermediate verification & What quantity the tube covers \\
Label denoising & Better synthetic labels & Coverage after optimization \\
PURE / shaping & Better credit assignment or outcome invariance & Tube benefits
with aggregation held fixed \\
Ensembles / AdvPO & Pessimistic reward optimization & Shared bias and joint
process uncertainty \\
Quantile / OT PRMs & Prefix distributions and bounds & Selection coverage and
completer dependence \\
Conformal / JOMI & Guarantees for specified sampling rules & Feasible PRM
adaptation and adaptive RL assumptions \\
Correlated-proxy robustness & Structured worst-case rewards & Construction from
process supervision \\
\bottomrule
\end{longtable}

\subsubsection{A defensible candidate contribution}
\emph{Optimization-aware uncertainty for process rewards}: compare pointwise
intervals with structured or problem-level tubes; measure coverage after
selection; determine whether coverage improvements reduce hacking at matched
optimization pressure. Novelty remains provisional. An empirical paper could
expose systematic post-selection undercoverage and a practical repair. A theory
paper could characterize safe improvement and failure of marginal intervals.
A combined paper would connect a restricted theorem to measured violations in
real PRMs. No publication outcome is assumed.

